\chapter{من المورّثة إلى البروتين}\label{ch:g11:gene-expression}

سنة 1961 غذّى كيميائي حيوي مستخلصًا خاليًا من \omterm{def:g10:cells-common-unit:cell}{الخلايا} من بكتيريا
بجزيء \omterm{def:g11:gene-expression:transcription}{RNA} اصطناعي مصنوع من الحرف U وحده، مكررًا. فصنع المستخلص بروتينًا
مصنوعًا من \omterm{def:g10:chemistry-of-life:families}{حمض أميني} واحد لا غير، هو الفينيل ألانين، مكررًا. وبتلك
التجربة الواحدة قرئت أول كلمة من \omterm{def:g11:gene-expression:code}{الشفرة الوراثية}: UUU تعني الفينيل
ألانين. وفي غضون خمس سنوات عرفت الكلمات الأربع والستون كلها، فتبين
أنها الكلمات نفسها في بكتيريا ونبتة قمح وإنسان. ويتابع هذا الفصل
الطريق من متتالية نوكليوتيدات \omterm{def:g10:universal-dna:gene}{المورّثة} إلى متتالية \omterm{def:g10:chemistry-of-life:families}{الحموض الأمينية} في
\omterm{def:g11:gene-expression:protein}{البروتين} --- وهو الطريق الذي تجري عليه كل صفة في
\cref{ch:g11:mutations}.

\section{المورّثات تصنع البروتينات}

\begin{proposition}[مورّثة واحدة، بروتين واحد]\label{prop:g11:gene-expression:onegene}
تعبَّر \omterm{def:g10:universal-dna:gene}{المورّثة} حين تستعمل \omterm{def:g10:cells-common-unit:cell}{الخلية} متتاليتها لصنع \omterm{def:g11:gene-expression:protein}{البروتين} الذي ترمز
له. ومعظم الصفات تتوقف على \omterm{def:g10:chemistry-of-life:families}{البروتينات} --- إنزيمات وألياف بنيوية ونواقل
ومستقبلات --- \omterm{def:g11:mutations:mutation}{والطفرة} تغير صفة بتغييرها \omterm{def:g11:gene-expression:protein}{البروتين} الذي تنتجه مورّثتها.
\end{proposition}
\begin{proof}[شواهد]
شعّع بيدل وتيتوم (1941) أبواغ عفن الخبز وجمعا الطافرات التي لم تعد
تستطيع النمو على وسط أدنى ما لم تضف مادة واحدة بعينها، \omterm{def:g10:chemistry-of-life:families}{كالحمض الأميني}
أرجينين. وكانت كل طافرة تنقصها إنزيمة واحدة من سلسلة التفاعلات التي
تصنع تلك المادة؛ وكان كل عيب يورث \omterm{def:g10:universal-dna:gene}{كمورّثة} واحدة؛ وكانت الطافرات
المحصورة عند خطوات مختلفة من السلسلة نفسها تحمل \omterm{def:g11:mutations:mutation}{طفرات} في \omterm{def:g10:universal-dna:gene}{مورّثات}
مختلفة. \omterm{def:g10:universal-dna:gene}{فمورّثة} واحدة، \omterm{def:g10:cell-metabolism:metabolism}{وإنزيم} واحد --- وكما عمّمت أعمال لاحقة، \omterm{def:g10:universal-dna:gene}{مورّثة}
واحدة، وسلسلة عديد ببتيد واحدة.
\end{proof}

\begin{definition}[البروتين ومتتالية الحموض الأمينية]\label{def:g11:gene-expression:protein}
\emph{البروتين}\index{بروتين} سلسلة من \omterm{def:g10:chemistry-of-life:families}{الحموض الأمينية}، طولها من بضع
عشرات إلى عدة آلاف، مأخوذة من مجموعة من عشرين صنفًا.
و\emph{المتتالية} --- أي ترتيب \omterm{def:g10:chemistry-of-life:families}{الحموض الأمينية} --- هي ما تحدده
\omterm{def:g10:universal-dna:gene}{المورّثة}؛ وحالما تصنع السلسلة تنطوي على شكل ثلاثي الأبعاد محدد تحدده
تلك المتتالية، والشكل يحدد الوظيفة. فغيّر حمضًا أمينيًا واحدًا فقد
يتغير الطي، وتتغير الوظيفة.
\end{definition}

\section{الاستنساخ: المورّثة تنسخ إلى RNA}

\begin{definition}[RNA الرسول والاستنساخ]\label{def:g11:gene-expression:transcription}
\emph{RNA}\index{RNA} (\omterm{def:g10:chemistry-of-life:families}{الحمض النووي} الريبي) سلسلة أحادية الشريط من
نوكليوتيدات كنوكليوتيدات \omterm{def:g10:universal-dna:information}{DNA}، إلا أن سكرها ريبوز وأن القاعدة يوراسيل
(U) تحل محل التيمين (T). و\emph{الاستنساخ}\index{استنساخ} هو نسخ
متتالية \omterm{def:g10:universal-dna:gene}{المورّثة} إلى جزيء من \emph{RNA الرسول}\index{RNA الرسول}
(mRNA): إذ يفتح \omterm{def:g10:cell-metabolism:metabolism}{الإنزيم} \emph{بوليميراز RNA} \omterm{prop:g10:universal-dna:helix}{اللولب المزدوج} على طول
\omterm{def:g10:universal-dna:gene}{المورّثة}، ويقرأ أحد الشريطين (\emph{الشريط الأصل}) ويركّب RNA المتمم
له، U في مواجهة A، وA في مواجهة T، وG في مواجهة C، وC في مواجهة G.
ولجزيء mRNA متتالية شريط \omterm{def:g10:universal-dna:gene}{المورّثة} الآخر، مع U بدل T؛ وهو يغادر \omterm{def:g10:cells-common-unit:organelle}{النواة} إلى
\omterm{def:g10:cells-common-unit:cell}{السيتوبلازم}.
\end{definition}

\begin{omfigure}
\begin{tikzpicture}[font=\small, >=Stealth, line cap=round]
  % DNA strands with bubble
  \draw[very thick, omDef] (0,0.5) -- (2.5,0.5) .. controls (3.2,0.5) and (3.4,1.3) .. (4.3,1.3) -- (6.7,1.3) .. controls (7.6,1.3) and (7.8,0.5) .. (8.5,0.5) -- (11,0.5);
  \draw[very thick, omDef] (0,-0.5) -- (2.5,-0.5) .. controls (3.2,-0.5) and (3.4,-1.3) .. (4.3,-1.3) -- (6.7,-1.3) .. controls (7.6,-1.3) and (7.8,-0.5) .. (8.5,-0.5) -- (11,-0.5);
  \foreach \x in {0.3,0.7,...,2.3} \draw[gray] (\x,0.5) -- (\x,-0.5);
  \foreach \x in {8.7,9.1,...,10.9} \draw[gray] (\x,0.5) -- (\x,-0.5);
  % template strand letters (bottom strand in bubble)
  \foreach \x/\l in {4.5/T, 4.9/A, 5.3/C, 5.7/G, 6.1/G, 6.5/A} \node[font=\scriptsize, omDef, anchor=north] at (\x,-1.32) {\l};
  % mRNA
  \draw[very thick, omThm] (4.4,-0.9) -- (6.6,-0.9) .. controls (7.4,-0.9) and (7.2,0.1) .. (8.0,0.6) .. controls (8.6,1.0) and (9.4,1.4) .. (10.6,1.7);
  \foreach \x/\l in {4.5/A, 4.9/U, 5.3/G, 5.7/C, 6.1/C, 6.5/U} \node[font=\scriptsize, omThm, anchor=south] at (\x,-0.85) {\l};
  \foreach \x in {4.5,4.9,...,6.5} \draw[gray, densely dotted] (\x,-0.95) -- (\x,-1.28);
  % polymerase
  \draw[thick, fill=omProp!30, rounded corners=8pt] (6.95,-1.6) rectangle (8.65,0.2);
  \node[align=center] at (7.8,-0.7) {بوليميراز\\ RNA};
  \draw[->, thick] (8.7,-1.9) -- (9.8,-1.9) node[right] {يتقدم على طول المورّثة};
  \node[anchor=east, omDef, align=right] at (-0.15,0) {DNA\\ المورّثة};
  \node[anchor=west, omThm, align=left] at (10.7,1.7) {mRNA، وهو\\ يتحرر};
  \node[anchor=north, omDef] at (5.5,-1.75) {الشريط الأصل};
\end{tikzpicture}

{\small \omterm{def:g11:gene-expression:transcription}{الاستنساخ}. بوليميراز \omterm{def:g11:gene-expression:transcription}{RNA} يفتح اللولب، ويقرأ الشريط الأصل،
ويبني \omterm{def:g11:gene-expression:transcription}{RNA} رسولًا متممًا له (U في مكان T). ووراء \omterm{def:g10:cell-metabolism:metabolism}{الإنزيم} ينغلق اللولب
من جديد؛ وينسلخ \omterm{def:g11:gene-expression:transcription}{RNA}.}
\end{omfigure}

\begin{example}[نسخة مستنسخة]\label{ex:g11:gene-expression:transcript}
الشريط الأصل \texttt{3'-TAC GGA CTT ATC-5'} يعطي mRNA
\texttt{5'-AUG CCU GAA UAG-3'}. وشريط \omterm{def:g10:universal-dna:gene}{المورّثة} الآخر يقرأ
\texttt{5'-ATG CCT GAA TAG-3'}: فجزيء mRNA هو متتالية ذلك الشريط، مع U بدل T، ولذلك تكتب متتاليات \omterm{def:g10:universal-dna:gene}{المورّثات} اصطلاحًا بذلك الشريط، وهو \emph{الشريط
المرمِّز}.
\end{example}

\section{الشفرة الوراثية}

\begin{definition}[الرامزة والشفرة الوراثية]\label{def:g11:gene-expression:code}
يقرأ mRNA في زمر متتابعة من ثلاثة نوكليوتيدات، هي
\emph{الرامزات}\index{رامزة}، من نقطة بدء محددة. و\emph{الشفرة
الوراثية}\index{شفرة وراثية} هي التقابل بين الرامزات الأربع والستين
الممكنة \omterm{def:g10:chemistry-of-life:families}{والحموض الأمينية} العشرين: فإحدى وستون رامزة تحدد حمضًا
أمينيًا، وثلاث (UAA وUAG وUGA) إشارات \emph{توقف} تنهي السلسلة،
وAUG، وهي تحدد الميثيونين، تعمل أيضًا رامزةَ \emph{بدء}. والشفرة
\emph{متعددة} --- فلمعظم \omterm{def:g10:chemistry-of-life:families}{الحموض الأمينية} عدة رامزات --- وهي
\emph{كونية}: فالجدول نفسه يخدم كل كائن معروف، مع متغايرات طفيفة نادرة.
\end{definition}

\begin{omfigure}
\begin{tabular}{c|cccc|c}
\multicolumn{1}{c}{} & \multicolumn{4}{c}{\small الحرف الثاني} & \multicolumn{1}{c}{} \\
{\small الأول} & U & C & A & G & {\small الثالث} \\ \hline
U & UUU Phe & UCU Ser & UAU Tyr & UGU Cys & U \\
 & UUC Phe & UCC Ser & UAC Tyr & UGC Cys & C \\
 & UUA Leu & UCA Ser & UAA \textbf{توقف} & UGA \textbf{توقف} & A \\
 & UUG Leu & UCG Ser & UAG \textbf{توقف} & UGG Trp & G \\ \hline
C & CUU Leu & CCU Pro & CAU His & CGU Arg & U \\
 & CUC Leu & CCC Pro & CAC His & CGC Arg & C \\
 & CUA Leu & CCA Pro & CAA Gln & CGA Arg & A \\
 & CUG Leu & CCG Pro & CAG Gln & CGG Arg & G \\ \hline
A & AUU Ile & ACU Thr & AAU Asn & AGU Ser & U \\
 & AUC Ile & ACC Thr & AAC Asn & AGC Ser & C \\
 & AUA Ile & ACA Thr & AAA Lys & AGA Arg & A \\
 & AUG \textbf{Met} & ACG Thr & AAG Lys & AGG Arg & G \\ \hline
G & GUU Val & GCU Ala & GAU Asp & GGU Gly & U \\
 & GUC Val & GCC Ala & GAC Asp & GGC Gly & C \\
 & GUA Val & GCA Ala & GAA Glu & GGA Gly & A \\
 & GUG Val & GCG Ala & GAG Glu & GGG Gly & G \\
\end{tabular}

{\small \omterm{def:g11:gene-expression:code}{الشفرة الوراثية}، مقروءةً على mRNA. تجد كل \omterm{def:g11:gene-expression:code}{رامزة} بحرفها الأول
(السطر) وحرفها الثاني (العمود) وحرفها الثالث (السطر داخل الكتلة). وAUG
هي الميثيونين وإشارة البدء معًا؛ وثلاث رامزات إشارات توقف.}
\end{omfigure}

\begin{proposition}[كيف قرئت الشفرة]\label{prop:g11:gene-expression:readcode}
أرسي معنى كل \omterm{def:g11:gene-expression:code}{رامزة} تجريبيًا.
\end{proposition}
\begin{proof}[شواهد]
أضاف نيرنبرغ وماتاي (1961) \omterm{def:g11:gene-expression:transcription}{RNA} اصطناعيًا من \omterm{def:g10:universal-dna:nucleotide}{نوكليوتيد} واحد مكرر إلى
مستخلص خال من \omterm{def:g10:cells-common-unit:cell}{الخلايا} يحوي ريبوزومات وحموضًا أمينية وإنزيمات تركيب
\omterm{def:g11:gene-expression:protein}{البروتين}: فأعطى متعدد U سلسلة من الفينيل ألانين، ومتعدد C سلسلة من
البرولين، ومتعدد A سلسلة من الليزين. ثم أسندت \omterm{def:g11:gene-expression:transcription}{RNA} اصطناعية من أزواج
وثلاثيات مكررة، ثم ارتباط الثلاثيات المفردة \omterm{def:g10:cells-common-unit:organelle}{بالريبوزومات}، الرامزاتِ
الباقية بحلول 1966. وقد استنتج طول الأحرف الثلاثة من الحساب (فحرفان لا
يعطيان إلا 16 تركيبًا، وهي قليلة جدًا لعشرين حمضًا أمينيًا؛ وثلاثة
تعطي 64) وأكدته \omterm{def:g11:mutations:mutation}{الطفرات}: فإقحام \omterm{def:g10:universal-dna:nucleotide}{نوكليوتيد} أو نوكليوتيدين في \omterm{def:g10:universal-dna:gene}{مورّثة}
يهدم بروتينها، وإقحام ثلاثة يعيد بروتينًا شبه طبيعي.
\end{proof}

\begin{method}[ترجمة متتالية]\label{met:g11:gene-expression:translate}
\begin{enumerate}
  \item اكتب الشريط المرمِّز (أو استنسخ الشريط الأصل): استبدل T بالحرف U
        لتحصل على mRNA.
  \item ابحث عن أول AUG: فهي البدء، ومنها يثبَّت إطار القراءة.
  \item قطّع mRNA إلى ثلاثيات متتابعة من AUG واقرأ كل ثلاثي في
        الجدول.
  \item قف عند أول \omterm{def:g11:gene-expression:code}{رامزة} توقف؛ \omterm{def:g10:chemistry-of-life:families}{والحموض الأمينية} المسرودة، بالترتيب،
        هي متتالية \omterm{def:g11:gene-expression:protein}{البروتين}.
  \item وللطافرة، أعد ذلك وقارن: \omterm{def:g11:gene-expression:protein}{البروتين} نفسه (صامتة)، أو \omterm{def:g10:chemistry-of-life:families}{حمض أميني}
        واحد تغير (خاطئة المعنى)، أو توقف سابق لأوانه (عديمة المعنى)،
        أو كل شيء تغير بعد \omterm{def:g11:mutations:mutation}{الطفرة} (إزاحة إطار).
\end{enumerate}
\end{method}

\begin{example}[أربع طفرات، أربع نتائج]\label{ex:g11:gene-expression:outcomes}
الشريط المرمِّز \texttt{ATG CCT GAA TTC TAG}: وmRNA
\texttt{AUG CCU GAA UUC UAG}، \omterm{def:g11:gene-expression:protein}{والبروتين} Met--Pro--Glu--Phe.
\begin{itemize}
  \item \texttt{ATG CC\textbf{C} GAA TTC TAG}: \omterm{def:g11:gene-expression:code}{فالرامزة} CCC لا تزال Pro ---
        وهي \emph{صامتة}.
  \item \texttt{ATG CCT G\textbf{C}A TTC TAG}: \omterm{def:g11:gene-expression:code}{فالرامزة} GCA هي Ala ---
        أي Met--Pro--Ala--Phe، وهي \emph{خاطئة المعنى}.
  \item \texttt{ATG CCT \textbf{T}AA TTC TAG}: \omterm{def:g11:gene-expression:code}{فالرامزة} UAA توقف ---
        أي Met--Pro، وهي \emph{عديمة المعنى}، أي \omterm{def:g11:gene-expression:protein}{بروتين} مبتور.
  \item \texttt{ATG CC\textbf{A}T GAA TTC TAG}: حرف A مقحم يزيح
        الإطار: \texttt{AUG CCA UGA}\dots{} أي Met--Pro--توقف ---
        وهي \emph{إزاحة إطار}.
\end{itemize}
\end{example}

\section{الترجمة: الرسالة تقرأ}

\begin{proposition}[الترجمة]\label{prop:g11:gene-expression:translation}
تجري \emph{الترجمة}\index{ترجمة} على \emph{\omterm{def:g10:cells-common-unit:organelle}{الريبوزومات}}، في
\omterm{def:g10:cells-common-unit:cell}{السيتوبلازم}. فيرتبط \omterm{def:g10:cells-common-unit:organelle}{الريبوزوم} بجزيء mRNA عند \omterm{def:g11:gene-expression:code}{رامزة} البدء ويتحرك عليه رامزةً
\omterm{def:g11:gene-expression:code}{رامزة}. وتقابل كل رامزةَ \emph{\omterm{def:g11:gene-expression:transcription}{RNA} ناقل} (tRNA)، وهو \omterm{def:g11:gene-expression:transcription}{RNA} صغير يحمل، في
أحد طرفيه، \emph{\omterm{def:g11:gene-expression:code}{الرامزة} المضادة} الثلاثية \omterm{def:g10:universal-dna:nucleotide}{النوكليوتيدات} المتممة
\omterm{def:g11:gene-expression:code}{للرامزة}، وفي الطرف الآخر \omterm{def:g10:chemistry-of-life:families}{الحمض الأميني} الموافق. ويصل \omterm{def:g10:cells-common-unit:organelle}{الريبوزوم} كل \omterm{def:g10:chemistry-of-life:families}{حمض
أميني} يجلب إلى السلسلة النامية ويطلق tRNA الفارغ؛ وعند \omterm{def:g11:gene-expression:code}{رامزة} توقف يطلق
السلسلة التامة، فتنطوي. وعدة ريبوزومات تقرأ mRNA واحدًا تباعًا، وmRNA
واحد يعطي مئات النسخ من \omterm{def:g11:gene-expression:protein}{البروتين} قبل أن يهدم.
\end{proposition}
\admitted

\begin{omfigure}
\begin{tikzpicture}[font=\small, >=Stealth, line cap=round]
  % mRNA
  \draw[very thick, omThm] (0,0) -- (11,0);
  \foreach \x/\l in {0.4/A,0.8/U,1.2/G, 1.7/C,2.1/C,2.5/U, 3.0/G,3.4/A,3.8/A, 4.3/U,4.7/U,5.1/C, 5.6/G,6.0/G,6.4/U, 6.9/U,7.3/A,7.7/G}
    \node[font=\scriptsize, omThm, anchor=north] at (\x,-0.05) {\l};
  \node[anchor=east, omThm] at (-0.15,0) {mRNA};
  % ribosome (two lobes) over codons 3 and 4
  \draw[thick, fill=omProp!25, rounded corners=10pt] (2.7,0.15) rectangle (5.5,1.4);
  \draw[thick, fill=omProp!35, rounded corners=12pt] (2.5,1.4) .. controls (2.5,3.0) and (5.7,3.0) .. (5.7,1.4) -- cycle;
  \node at (4.1,2.1) {ريبوزوم};
  % tRNAs
  \draw[very thick, omMeth!80!black] (3.4,0.55) -- (3.4,1.35);
  \node[font=\scriptsize, omMeth!80!black, anchor=south] at (3.4,0.15) {CUU};
  \draw[very thick, omMeth!80!black] (4.7,0.55) -- (4.7,1.35);
  \node[font=\scriptsize, omMeth!80!black, anchor=south] at (4.7,0.15) {AAG};
  % amino acids chain
  \foreach \i/\x in {0/3.4, 1/4.7} \fill[omDef!70] (\x,1.5) circle (4pt);
  \foreach \x in {2.6, 1.9, 1.2} \fill[omDef!70] (\x,1.9) circle (4pt);
  \draw[thick, omDef!70] (1.2,1.9) -- (2.6,1.9) -- (3.4,1.5) -- (4.7,1.5);
  \node[anchor=east, omDef!80!black, align=right] at (0.9,1.9) {البروتين\\ النامي};
  % incoming tRNA
  \draw[very thick, omMeth!80!black] (7.0,1.3) -- (7.0,2.1);
  \fill[omDef!70] (7.0,2.25) circle (4pt);
  \node[font=\scriptsize, omMeth!80!black, anchor=west] at (7.15,1.35) {CCA};
  \node[anchor=west, align=left, font=\scriptsize] at (7.3,1.7) {tRNA التالي، ومعه\\ حمضه الأميني، يصل};
  \draw[->, thick] (7.0,1.1) -- (6.1,0.75);
  \draw[->, thick] (5.7,-0.9) -- (7.2,-0.9) node[right] {الريبوزوم يتقدم};
  \node[anchor=north, font=\scriptsize] at (4.05,-0.45) {رامزة};
  \draw (2.85,-0.35) -- (2.85,-0.45) -- (3.95,-0.45) -- (3.95,-0.35);
\end{tikzpicture}

{\small \omterm{prop:g11:gene-expression:translation}{الترجمة}. \omterm{def:g10:cells-common-unit:organelle}{الريبوزوم} يمسك رامزتين؛ وناقل tRNA تطابق رامزته
المضادة كل \omterm{def:g11:gene-expression:code}{رامزة} يجلب حمضه الأميني، فتنقل السلسلة إلى القادم الجديد،
ويخطو \omterm{def:g10:cells-common-unit:organelle}{الريبوزوم} رامزةً واحدة إلى الأمام. وعند \omterm{def:g11:gene-expression:code}{رامزة} توقف تطلق
السلسلة.}
\end{omfigure}

\begin{example}[السرعة والمردود]\label{ex:g11:gene-expression:speed}
يضيف \omterm{def:g10:cells-common-unit:organelle}{الريبوزوم} 5 إلى 20 حمضًا أمينيًا في الثانية: \omterm{def:g11:gene-expression:protein}{فبروتين} من 300 \omterm{def:g10:chemistry-of-life:families}{حمض
أميني} يستغرق دون دقيقة. \omterm{def:g10:cells-common-unit:organelle}{والريبوزومات} تتبع بعضها بعضًا على mRNA على
مسافة نحو 80 نوكليوتيدًا، فرسالة من 900 \omterm{def:g10:universal-dna:nucleotide}{نوكليوتيد} تحمل عشرة منها في آن
واحد؛ وعلى مدى حياة تدوم بضع ساعات ينتج mRNA واحد ألف جزيء بروتيني.
\omterm{def:g10:cells-common-unit:cell}{وخلية} \emph{E.~coli} تصنع نحو $\num{20000}$ جزيء بروتيني في الثانية؛
وسليف الكرية الحمراء نحو $\num{5e8}$ جزيء هيموغلوبين في حياته.
\end{example}

\section{مورّثة واحدة، رسائل عدة}

\begin{proposition}[نضج الرسالة في حقيقيات النواة]\label{prop:g11:gene-expression:splicing}
في \omterm{def:g10:cells-common-unit:cell}{الخلايا} الحقيقية \omterm{def:g10:cells-common-unit:organelle}{النواة} تتخلل متتالية \omterm{def:g10:universal-dna:gene}{المورّثة}
\emph{إنترونات}\index{إنترون}، وهي قطع تستنسخ ثم تقطع من \omterm{def:g11:gene-expression:transcription}{RNA} قبل أن
يغادر \omterm{def:g10:cells-common-unit:organelle}{النواة}؛ أما القطع التي تحفظ وتوصل طرفًا بطرف، وهي
\emph{الإكسونات}، فتكوّن mRNA الناضج. ويمكن قطع \omterm{def:g10:universal-dna:gene}{مورّثات} كثيرة بأكثر من
طريقة، فتحفظ مجموعات مختلفة من الإكسونات (\emph{التضفير البديل})، حتى
تعطي \omterm{def:g10:universal-dna:gene}{مورّثة} واحدة عدة رسائل mRNA متمايزة وعدة بروتينات قريبة: \omterm{def:g10:universal-dna:gene}{فمورّثات}
الإنسان $\num{20000}$ ترمز لما يزيد كثيرًا على $\num{100000}$ \omterm{def:g11:gene-expression:protein}{بروتين}.
\end{proposition}
\admitted

\begin{omfigure}
\begin{tikzpicture}[font=\small, >=Stealth]
  \tikzset{ex/.style={draw, fill=omDef!40, minimum height=0.45cm, minimum width=1.0cm, inner sep=1pt},
           intr/.style={draw, fill=gray!20, minimum height=0.45cm, minimum width=0.7cm, inner sep=1pt}}
  \node[anchor=east] at (-0.2,0) {المورّثة (DNA)};
  \node[ex] at (0.5,0) {1}; \node[intr] at (1.35,0) {}; \node[ex] at (2.2,0) {2}; \node[intr] at (3.05,0) {}; \node[ex] at (3.9,0) {3}; \node[intr] at (4.75,0) {}; \node[ex] at (5.6,0) {4};
  \node[anchor=west, font=\scriptsize] at (6.2,0) {إكسونات (بالأزرق)، إنترونات (بالرمادي)};
  \draw[->, thick] (3.0,-0.35) -- (3.0,-0.95) node[midway, right] {استنساخ};
  \node[anchor=east] at (-0.2,-1.3) {mRNA أولي};
  \node[ex, fill=omThm!30] at (0.5,-1.3) {1}; \node[intr] at (1.35,-1.3) {}; \node[ex, fill=omThm!30] at (2.2,-1.3) {2}; \node[intr] at (3.05,-1.3) {}; \node[ex, fill=omThm!30] at (3.9,-1.3) {3}; \node[intr] at (4.75,-1.3) {}; \node[ex, fill=omThm!30] at (5.6,-1.3) {4};
  \draw[->, thick] (2.0,-1.65) -- (1.0,-2.35) node[midway, left] {تضفير};
  \draw[->, thick] (4.0,-1.65) -- (6.0,-2.35) node[midway, right] {تضفير بديل};
  \node[anchor=east] at (-0.6,-2.7) {mRNA A};
  \node[ex, fill=omThm!30] at (0.0,-2.7) {1}; \node[ex, fill=omThm!30] at (1.0,-2.7) {2}; \node[ex, fill=omThm!30] at (2.0,-2.7) {3}; \node[ex, fill=omThm!30] at (3.0,-2.7) {4};
  \node[anchor=west] at (3.6,-2.7) {البروتين A};
  \node[anchor=east] at (5.4,-2.7) {};
  \node[ex, fill=omThm!30] at (6.0,-2.7) {1}; \node[ex, fill=omThm!30] at (7.0,-2.7) {2}; \node[ex, fill=omThm!30] at (8.0,-2.7) {4};
  \node[anchor=west] at (8.6,-2.7) {البروتين B};
\end{tikzpicture}

{\small \omterm{def:g10:universal-dna:gene}{المورّثة} الحقيقية \omterm{def:g10:cells-common-unit:organelle}{النواة} تستنسخ كاملة، ثم تزال \omterm{prop:g11:gene-expression:splicing}{الإنترونات}.
وحفظ الإكسونات كلها، أو ترك واحد منها، يعطي رسالتي mRNA وبروتينين من
\omterm{def:g10:universal-dna:gene}{المورّثة} نفسها.}
\end{omfigure}

\begin{remark}[جريان المعلومة]\label{rem:g11:gene-expression:flow}
\omterm{def:g10:universal-dna:information}{DNA} يستنسخ إلى \omterm{def:g11:gene-expression:transcription}{RNA}، \omterm{def:g11:gene-expression:transcription}{وRNA} يترجم إلى \omterm{def:g11:gene-expression:protein}{بروتين}، \omterm{def:g11:gene-expression:protein}{والبروتين} يصنع الصفة:
فالمعلومة تجري في اتجاه واحد. \omterm{def:g11:gene-expression:protein}{والبروتين} المتغير لا يعيد كتابة
\omterm{def:g10:universal-dna:gene}{المورّثة} أبدًا، ولذلك لا يورث عمر من \omterm{def:g10:sport-and-health:training}{التدريب} ولا من حروق الشمس، ولذلك
لا يغير ما يتلقاه الجيل التالي إلا \omterm{def:g11:mutations:mutation}{طفرات} \omterm{def:g10:universal-dna:information}{DNA}
(\cref{ch:g11:mutations}). والخطوات الثلاث نفسها تجري في كل \omterm{def:g10:cells-common-unit:cell}{خلية} من كل
كائن، وبالشفرة نفسها: أي كونية \cref{ch:g10:universal-dna}، نازلةً
الآن إلى آخر كلمة.
\end{remark}

\section{تمارين}

\begin{exercise}[$\star$]\label{exo:g11:gene-expression:1}
أعطِ ثلاثة فروق بين \omterm{def:g10:universal-dna:information}{DNA} \omterm{def:g11:gene-expression:transcription}{وRNA}.
\end{exercise}

\begin{exercise}[$\star$]\label{exo:g11:gene-expression:2}
استنسخ الشريط الأصل \texttt{3'-TAC CGA AAA ATT-5'} إلى mRNA.
\end{exercise}

\begin{exercise}[$\star$]\label{exo:g11:gene-expression:3}
ترجم mRNA \texttt{AUG GGC AAA UGU UAA} بجدول الشفرة.
\end{exercise}

\begin{exercise}[$\star$]\label{exo:g11:gene-expression:4}
ما دورا \omterm{def:g10:cells-common-unit:organelle}{الريبوزوم} \omterm{def:g11:gene-expression:transcription}{وRNA} الناقلة في \omterm{prop:g11:gene-expression:translation}{الترجمة}؟
\end{exercise}

\begin{exercise}[$\star$]\label{exo:g11:gene-expression:5}
لماذا يجب أن تستعمل الشفرة ثلاثة نوكليوتيدات على الأقل لكل \omterm{def:g10:chemistry-of-life:families}{حمض أميني}؟
\end{exercise}

\begin{exercise}[$\star\star$]\label{exo:g11:gene-expression:6}
\omterm{def:g11:gene-expression:protein}{بروتين} فيه 412 حمضًا أمينيًا. ما أقل طول لجزئه المرمِّز من mRNA، مع
\omterm{def:g11:gene-expression:code}{رامزة} التوقف، وما أقل طول \omterm{def:g10:universal-dna:gene}{للمورّثة} الموافقة بأزواج القواعد؟
\end{exercise}

\begin{exercise}[$\star\star$]\label{exo:g11:gene-expression:7}
الشريط المرمِّز \texttt{ATG AAA GGC TGG TAA} يطفر إلى
\texttt{ATG AAA GGC TGA TAA}. أعطِ البروتينين وصنّف \omterm{def:g11:mutations:mutation}{الطفرة}.
\end{exercise}

\begin{exercise}[$\star\star$]\label{exo:g11:gene-expression:8}
المتتالية الابتدائية نفسها، تطفر إلى \texttt{ATG AAG GGC TGG TAA} وإلى
\texttt{ATG AAA GGG CTG GTA A}. صنّف كلًا منهما وأعطِ البروتينين.
\end{exercise}

\begin{exercise}[$\star\star$]\label{exo:g11:gene-expression:9}
فسّر لماذا يكون الاستبدال في الموضع الثالث من \omterm{def:g11:gene-expression:code}{الرامزة} صامتًا غالبًا،
باستعمال الجدول.
\end{exercise}

\begin{exercise}[$\star\star$]\label{exo:g11:gene-expression:10}
\omterm{def:g11:gene-expression:transcription}{RNA} من متعدد UC (\texttt{UCUCUCUC}\dots) يعطي بروتينًا يتناوب فيه
السيرين واللوسين. بيّن أن ذلك متسق مع شفرة من ثلاثة أحرف واستعمله
لإسناد رامزتين.
\end{exercise}

\begin{exercise}[$\star\star$]\label{exo:g11:gene-expression:11}
\omterm{def:g10:universal-dna:gene}{مورّثة} بشرية من \num{30000} زوج قواعد تنتج بروتينًا من 500 \omterm{def:g10:chemistry-of-life:families}{حمض أميني}.
فسّر التباين.
\end{exercise}

\begin{exercise}[$\star\star\star$]\label{exo:g11:gene-expression:12}
دواء يحصر ريبوزومات البكتيريا ولا يحصر ريبوزومات الإنسان. فسّر لماذا
يستطيع أن يكون مضادًا حيويًا، وما الذي يعنيه وجوده عن \omterm{def:g10:cells-common-unit:organelle}{الريبوزومات} في
المجموعتين.
\end{exercise}

\begin{exercise}[$\star\star\star$]\label{exo:g11:gene-expression:13}
فسّر كيف استطاع فأر أن يقرأ \omterm{def:g10:universal-dna:gene}{مورّثة} قنديل البحر في
\cref{ch:g10:universal-dna}، بمفردات هذا الفصل، وماذا كان ليحدث لو
استعمل الفأر شفرة أخرى.
\end{exercise}

\begin{exercise}[$\star\star\star$]\label{exo:g11:gene-expression:14}
\omterm{def:g11:mutations:mutation}{طفرة} تغير tRNA الذي تقرأ رامزته المضادة UAG فيصير يحمل حمضًا أمينيًا
بدل أن يوقف. تنبأ بأثرها في بروتينات \omterm{def:g10:cells-common-unit:cell}{الخلية} عمومًا، وفي \omterm{def:g10:universal-dna:gene}{مورّثة} طافرة
تحمل UAG سابقة لأوانها.
\end{exercise}

\begin{exercise}[$\star\star\star$]\label{exo:g11:gene-expression:15}
ناقش لماذا تكون إزاحة الإطار قرب بداية \omterm{def:g10:universal-dna:gene}{المورّثة} أسوأ دائمًا تقريبًا من
إزاحتها قرب نهايتها، ولماذا تستطيع \omterm{def:g11:mutations:mutation}{الطفرة} الخاطئة المعنى أن تكون أي
شيء بين غير المؤذي والمميت.
\end{exercise}

\section{مسألة: مورّثة تقرأ على أربعة أوجه}

\begin{problem}[{مسألة نهاية الأسبوع --- مورّثة قصيرة واحدة تستنسخ،
وتترجم، وتطفر ثلاث مرات، ويقاس زمنها: من DNA إلى الحموض الأمينية،
والثواني الستون التي يستغرقها بروتين}]
\label{pb:g11:gene-expression:1}
الشريط المرمِّز \omterm{def:g10:universal-dna:gene}{لمورّثة} قصيرة يقرأ:
\begin{center}
\texttt{ATG GCT TGG AAA CCC GAA TAC GGT CAT TTC TGA}
\end{center}

\medskip
\noindent\textbf{الجزء الأول --- قراءة \omterm{def:g10:universal-dna:gene}{المورّثة}.}
\begin{enumerate}
  \item اكتب الشريط الأصل.
  \item اكتب mRNA المستنسخ من \omterm{def:g10:universal-dna:gene}{المورّثة}.
  \item ترجمه، رامزةً \omterm{def:g11:gene-expression:code}{رامزة}، وأعطِ متتالية \omterm{def:g10:chemistry-of-life:families}{الحموض الأمينية} في
        \omterm{def:g11:gene-expression:protein}{البروتين}.
  \item كم حمضًا أمينيًا يحوي \omterm{def:g11:gene-expression:protein}{البروتين}، وكم نوكليوتيدًا من mRNA يستعمل
        لتحديدها، مع \omterm{def:g11:gene-expression:code}{رامزة} التوقف؟
  \item أي \omterm{def:g11:gene-expression:code}{رامزة} بدأت القراءة، وما الذي يثبّت إطار القراءة لسائر
        \omterm{def:g11:gene-expression:code}{الرامزات}؟
\end{enumerate}

\medskip
\noindent\textbf{الجزء الثاني --- ثلاث طافرات.}
\begin{enumerate}[resume]
  \item الطافرة 1: \texttt{ATG GCT TGG AAA CCG GAA TAC GGT CAT TTC TGA}.
        أعطِ \omterm{def:g11:gene-expression:protein}{البروتين} وصنّف \omterm{def:g11:mutations:mutation}{الطفرة}.
  \item الطافرة 2: \texttt{ATG GCT TGG AAA CCC GAA TAC GGT CAT TTC TGA}
        مع استبدال \omterm{def:g11:gene-expression:code}{الرامزة} التاسعة \texttt{CAT} \omterm{def:g11:gene-expression:code}{بالرامزة} \texttt{CGT}.
        أعطِ \omterm{def:g11:gene-expression:protein}{البروتين} وصنّف.
  \item الطافرة 3: \texttt{ATG GCT TGA AAA CCC GAA TAC GGT CAT TTC TGA}.
        أعطِ \omterm{def:g11:gene-expression:protein}{البروتين} وصنّف. وأي استبدال واحد أنتجها؟
  \item الطافرة 4: حرف T مقحم بعد \omterm{def:g10:universal-dna:nucleotide}{النوكليوتيد} السادس:
        \texttt{ATG GCT TTG GAA ACC CGA ATA CGG TCA TTT CTG A}. أعطِ
        \omterm{def:g11:gene-expression:protein}{البروتين} وصنّف.
  \item رتّب الطافرات الأربع من الأقل إلى الأشد إخلالًا \omterm{def:g11:gene-expression:protein}{بالبروتين}
        وعلّل.
\end{enumerate}

\medskip
\noindent\textbf{الجزء الثالث --- لماذا كانت الشفرة كما هي.}
\begin{enumerate}[resume]
  \item كم \omterm{def:g11:gene-expression:code}{رامزة} تحدد اللوسين؟ والسيرين؟ والتربتوفان؟ وأي \omterm{def:g10:chemistry-of-life:families}{حمض أميني}
        أرجح أن يغيّره استبدال عشوائي في رامزته، وأيها أقلها ترجيحًا؟
  \item احسب نسبة الاستبدالات الصامتة من جميع الاستبدالات في الموضع
        الثالث من \omterm{def:g11:gene-expression:code}{رامزة} اللوسين \texttt{CUx}.
  \item فسّر لماذا لم تكن شفرة من حرفين لتنجح، ولماذا لا حاجة إلى شفرة
        من أربعة.
  \item متعدد U يعطي متعدد الفينيل ألانين، ومتعدد A يعطي متعدد
        الليزين. فأي مدخلات الجدول ترسيها هاتان التجربتان؟
  \item الشفرة واحدة في البكتيريا وفي الإنسان. اذكر النتيجة بالنسبة
        إلى \omterm{prop:g10:universal-dna:universal}{النقل المورّثي} والنتيجة بالنسبة إلى القرابة.
\end{enumerate}

\medskip
\noindent\textbf{الجزء الرابع --- توقيت المصنع.}
يضيف \omterm{def:g10:cells-common-unit:organelle}{الريبوزوم} 10 حموض أمينية في الثانية \omterm{def:g10:cells-common-unit:organelle}{والريبوزومات} تتبع بعضها بعضًا
على مسافة 80 نوكليوتيدًا على mRNA؛ ويعيش كل mRNA ساعتين.
\begin{enumerate}[resume]
  \item كم يستغرق \omterm{def:g10:cells-common-unit:organelle}{ريبوزوم} واحد لترجمة \omterm{def:g11:gene-expression:protein}{بروتين} هذه \omterm{def:g10:universal-dna:gene}{المورّثة}؟ \omterm{def:g11:gene-expression:protein}{وبروتين} من
        600 \omterm{def:g10:chemistry-of-life:families}{حمض أميني}؟
  \item كم ريبوزومًا يستطيع قراءة mRNA من 1800 \omterm{def:g10:universal-dna:nucleotide}{نوكليوتيد} في آن واحد؟
  \item إذا بدأ \omterm{def:g10:cells-common-unit:organelle}{ريبوزوم} جديد كل 8 ثوان، فكم نسخة من \omterm{def:g11:gene-expression:protein}{البروتين} يعطي mRNA
        واحد في حياته؟
  \item تحتاج بكتيريا إلى \num{2000} نسخة من \omterm{def:g10:cell-metabolism:metabolism}{إنزيم} في 10 دقائق. فكم
        رسالة mRNA من \omterm{def:g10:universal-dna:gene}{المورّثة}، تستنسخ في آن واحد، يلزم؟
  \item اذكر النتيجة: \omterm{def:g11:gene-expression:protein}{البروتين} الذي ترمز له \omterm{def:g10:universal-dna:gene}{المورّثة}، \omterm{def:g11:mutations:mutation}{والطفرة} التي
        ألغته من بين الأربع، والزمن الذي يحتاج إليه \omterm{def:g10:cells-common-unit:organelle}{ريبوزوم} واحد
        لبنائه.
\end{enumerate}
\end{problem}
